
% BEGIN bib
 
\par\medskip\noindent\textbf{References}\par
  
\par\noindent  [1]  Samson Abramsky.  No-cloning in categorical quantum mechanics.  Semantic Techniques in Quantum Computation, pages 1–28, 2009.  
 
\par\noindent  [2]  Leonard M. Adleman, Jonathan DeMarrais, and Ming-Deh A. Huang.  Quantum computability.  SIAM Journal on Computing, 26(5):1524–1540, 1997.  doi:10.1137/S0097539795293639.  
 
\par\noindent  [3]  Thorsten Altenkirch and Jonathan Grattage.  A functional quantum programming language.  In 20th Annual IEEE Symposium on Logic in Computer Science (LICS’ 05), pages 249–258, 2005.  doi:10.1109/LICS.2005.1.  
 
\par\noindent  [4]  John Baez.  Quantum quandaries: a category-theoretic perspective.  The structural foundations of quantum gravity, pages 240–265, 2006.  
 
\par\noindent  [5]  Adriano Barenco, Charles H. Bennett, Richard Cleve, David P. DiVincenzo, Norman Margolus, Peter Shor, Tycho Sleator, John A. Smolin, and Harald Weinfurter.  Elementary gates for quantum computation.  Phys. Rev. A, 52:3457–3467, November 1995.  doi:10.1103/PhysRevA.52.3457.  
 
\par\noindent  [6]  Ethan Bernstein and Umesh Vazirani.  Quantum complexity theory.  SIAM Journal on Computing, 26(5):1411–1473, 1997.  doi:10.1137/S0097539796300921.  
 
\par\noindent  [7]  Hans J. Briegel, David E. Browne, Wolfgang Dür, Robert Raussendorf, and Maarten Van den Nest.  Measurement-based quantum computation.  Nature Physics, 5(1):19–26, 2009.  doi:10.1038/nphys1157.  
 
\par\noindent  [8]  Andrew M. Childs and Nathan Wiebe.  Hamiltonian simulation using linear combinations of unitary operations.  Quantum Info. Comput., 12(11–12):901–924, November 2012.  doi:10.26421/QIC12.11-12-1.  
 
\par\noindent  [9]  Andrea Colledan and Ugo Dal Lago.  Circuit width estimation via effect typing and linear dependency.  In European Symposium on Programming, pages 3–30. Springer, 2024.  doi:10.1007/978-3-031-57267-8\_1.  
 
\par\noindent  [10]  Andrea Colledan and Ugo Dal Lago.  Flexible Type-Based Resource Estimation in Quantum Circuit Description Languages.  Proc. ACM Program. Lang., 9(POPL), January 2025.  doi:10.1145/3704883.  
 
\par\noindent  [11]  Ugo Dal Lago, Andrea Masini, and Margherita Zorzi.  Quantum implicit computational complexity.  Theoretical Computer Science, 411(2):377–409, 2010.  doi:10.1016/j.tcs.2009.07.045.  
 
\par\noindent  [12]  Vincent Danos and Elham Kashefi.  Determinism in the one-way model.  Physical Review A, 74(5):052310, 2006.  doi:10.1103/PhysRevA.74.052310.  
 
\par\noindent  [13]  Niel de Beaudrap, Xiaoning Bian, and Quanlong Wang.  Fast and Effective Techniques for T-Count Reduction via Spider Nest Identities.  In 15th Conference on the Theory of Quantum Computation, Communication and Cryptography (TQC 2020), pages 11:1–11:23. Schloss Dagstuhl – Leibniz-Zentrum für Informatik, 2020.  doi:10.4230/LIPICS.TQC.2020.11.  
 
\par\noindent  [14]  Alejandro Díaz-Caro, Mauricio Guillermo, Alexandre Miquel, and Benoît Valiron.  Realizability in the unitary sphere.  In 2019 34th Annual ACM/IEEE Symposium on Logic in Computer Science (LICS), pages 1–13. IEEE, 2019.  doi:10.1109/LICS.2019.8785834.  
 
\par\noindent  [15]  Peng Fu, Kohei Kishida, and Peter Selinger.  Linear dependent type theory for quantum programming languages: Extended abstract.  In Proceedings of the 35th Annual ACM/IEEE Symposium on Logic in Computer Science, LICS ’20, pages 440–453, New York, NY, USA, 2020. Association for Computing Machinery.  doi:10.1145/3373718.3394765.  
 
\par\noindent  [16]  David Gosset, Robin Kothari, and Kewen Wu.  Quantum state preparation with optimal T-count.  arXiv preprint arXiv:2411.04790, 2024.  doi:10.48550/arXiv.2411.04790.  
 
\par\noindent  [17]  Emmanuel Hainry, Romain Péchoux, and Mário Silva.  A programming language characterizing quantum polynomial time.  In Foundations of Software Science and Computation Structures, pages 156–175. Springer, 2023.  doi:10.1007/978-3-031-30829-1\_8.  
 
\par\noindent  [18]  Aram W. Harrow, Avinatan Hassidim, and Seth Lloyd.  Quantum Algorithm for Linear Systems of Equations.  Physical Review Letters, 103(15), October 2009.  doi:10.1103/physrevlett.103.150502.  
 
\par\noindent  [19]  Chris Heunen and Jamie Vicary.  Categories for Quantum Theory: an introduction.  Oxford University Press, 2019.  
 
\par\noindent  [20]  Aleks Kissinger and John van de Wetering.  Reducing the number of non-Clifford gates in quantum circuits.  Phys. Rev. A, 102:022406, August 2020.  doi:10.1103/PhysRevA.102.022406.  
 
\par\noindent  [21]  Emanuel Knill, Raymond Laflamme, and Gerard J. Milburn.  A scheme for efficient quantum computation with linear optics.  Nature, 409(6816):46–52, January 2001.  doi:10.1038/35051009.  
 
\par\noindent  [22]  Gui-Lu Long and Yang Sun.  Efficient scheme for initializing a quantum register with an arbitrary superposed state.  Phys. Rev. A, 64:014303, June 2001.  doi:10.1103/PhysRevA.64.014303.  
 
\par\noindent  [23]  Dmitri Maslov, Gerhard W. Dueck, D. Michael Miller, and Camille Negrevergne.  Quantum Circuit Simplification and Level Compaction.  IEEE Transactions on Computer-Aided Design of Integrated Circuits and Systems, 27(3):436–444, March 2008.  doi:10.1109/tcad.2007.911334.  
 
\par\noindent  [24]  Cristopher Moore and Martin Nilsson.  Parallel Quantum Computation and Quantum Codes.  SIAM Journal on Computing, 31(3):799–815, 2001.  doi:10.1137/S0097539799355053.  
 
\par\noindent  [25]  Yunseong Nam, Neil J. Ross, Yuan Su, Andrew M. Childs, and Dmitri Maslov.  Automated optimization of large quantum circuits with continuous parameters.  npj Quantum Information, 4(1):23, May 2018.  doi:10.1038/s41534-018-0072-4.  
 
\par\noindent  [26]  Michael A. Nielsen and Isaac L. Chuang.  Quantum Computation and Quantum Information: 10th Anniversary Edition.  Cambridge University Press, 2011.  
 
\par\noindent  [27]  Peter Selinger.  Towards a quantum programming language.  Mathematical Structures in Computer Science, 14(4):527–586, 2004.  doi:10.1017/S0960129504004256.  
 
\par\noindent  [28]  Xiaoming Sun, Guojing Tian, Shuai Yang, Pei Yuan, and Shengyu Zhang.  Asymptotically optimal circuit depth for quantum state preparation and general unitary synthesis.  IEEE Transactions on Computer-Aided Design of Integrated Circuits and Systems, 42(10):3301–3314, 2023.  doi:10.1109/TCAD.2023.3244885.  
 
\par\noindent  [29]  Tomoyuki Yamakami.  A Schematic Definition of Quantum Polynomial Time Computability.  J. Symb. Log., 85(4):1546–1587, 2020.  doi:10.1017/jsl.2020.45.  
 
\par\noindent  [30]  Tomoyuki Yamakami.  Expressing Power of Elementary Quantum Recursion Schemes for Quantum Logarithmic-Time Computability.  In Agata Ciabattoni, Elaine Pimentel, and Ruy J. G. B. de Queiroz, editors, Logic, Language, Information, and Computation, pages 88–104, Cham, 2022. Springer International Publishing.  doi:10.1007/978-3-031-15298-6\_6.  
 
\par\noindent  [31]  Andrew Chi-Chih Yao.  Quantum circuit complexity.  In Proceedings of 1993 IEEE 34th Annual Foundations of Computer Science, pages 352–361, 1993.  doi:10.1109/SFCS.1993.366852.  
 
\par\noindent  [32]  Charles Yuan and Michael Carbin.  Tower: Data Structures in Quantum Superposition.  Proceedings of the ACM on Programming Languages, 6(OOPSLA2):259–288, October 2022.  doi:10.1145/3563297.  
 
\par\noindent  [33]  Charles Yuan, Agnes Villanyi, and Michael Carbin.  Quantum control machine: The limits of control flow in quantum programming.  Proc. ACM Program. Lang., 8(OOPSLA1), April 2024.  doi:10.1145/3649811.  
 
\par\noindent  [34]  Xiao-Ming Zhang, Tongyang Li, and Xiao Yuan.  Quantum State Preparation with Optimal Circuit Depth: Implementations and Applications.  Phys. Rev. Lett., 129:230504, November 2022.  doi:10.1103/PhysRevLett.129.230504.  
 
\par\noindent  [35]  Zhicheng Zhang and Mingsheng Ying.  Quantum Register Machine: Efficient Implementation of Quantum Recursive Programs, 2024.  doi:10.48550/arXiv.2408.10054.  
  
% END bib
